\section{Zero-residual protection and good missing edges}
\label{sec:good-missing-residual}

Throughout the main argument, $D$ is nonempty, strongly connected, all-deficient,
and minimal under deletion of arbitrary nonempty sets of arcs.
No assumption $n=2\delta+3$ is made.
Write
\[
 g(x)=d_D^+(x)-d_D^{++}(x)>0,\qquad
 \eta(x)=\mu(x)-1+d_D^{--}(x)-d_D^{++}(x).
\]
The established pointwise residual theorem gives
\[
 \eta(x)\ge0,\qquad \sum_x\eta(x)=2M-n,
\]
where $M$ is the number of missing unordered pairs.

\subsection{A zero-residual source protects all of its far targets}

\begin{lemma}\label{lem:zero-residual-source-protection}
If $\eta(x)=0$ and $y\in U_D(x)$, then
\[
 y\to x,\qquad N_D^-(y)\subseteq N_D^-(x)\setminus\{y\}.
\]
\end{lemma}
\begin{proof}
Use the target-bundle sets
\[
 W=N_D^-(x),\quad T=V(D)\setminus(\{x\}\cup W),\quad
 P=\{v:x\in U_D(v)\},\quad R=T\setminus P=N_D^{--}(x).
\]
If $P$ is empty, $\eta(x)=2\mu(x)+g(x)-1=0$ forces
$\mu(x)=0$ and $g(x)=1$.
A far target $y$ then points to $x$, since $x$ is whole.
No outneighbor of $x$ points to $y$, so every predecessor of $y$
belongs to $N_D^-(x)\setminus\{y\}$.

Now let $P\ne\varnothing$ and put $C=\partial P$,
$X=\partial^2P$, $j=|R\setminus C|$ and
$s=|N_D^+(x)\cap(R\setminus C)|$.
Strong connectivity gives $C\ne\varnothing$.
The established sharpened target-bundle bounds are
\[
 s=0\Longrightarrow\eta(x)\ge2j,\qquad
 s>0\Longrightarrow\eta(x)\ge g(x)+2(j-s).
\]
Since $j\ge s$ and $g(x)>0$, zero residual forces $j=s=0$.
Thus $C=R$, $T=P\dot\cup C$, and
$|X|\le |C|-1$ by the external-boundary lemma.
Every second outneighbor of $x$ lies in
$\mathcal M(x)\cup X$.
Since $C\subseteq N_D^{--}(x)$, no vertex of $C$ points directly
to $x$. Thus $X\subseteq W$ and $\mathcal M(x)\subseteq T$, so these
sets are disjoint.  Zero residual gives
\[
 d_D^{++}(x)=\mu(x)+|C|-1.
\]
Equality of cardinalities therefore forces
\[
 N_D^{++}(x)=\mathcal M(x)\dot\cup X,\qquad |X|=|C|-1.
\]
Since $T=N_D^+(x)\dot\cup\mathcal M(x)$, this yields
$U_D(x)=W\setminus X$.
For $y\in U_D(x)$, there is no arc from $P$ to $y$ because
all out-arcs of $P$ remain in $T$.
There is no arc from $C$ to $y$, because such a head outside
$P\cup C$ would belong to $X$.
Also $x\not\to y$.  Hence $y\to x$ and all predecessors of $y$
are in $W\setminus\{y\}$, as required.
\end{proof}

\subsection{The published good-missing-edge transfer}

A missing pair $ab$ is \emph{good} in the sense of Ghazal if either
\[
 \forall w\notin\{a,b\}:\quad
 w\to a\ \Longrightarrow\ b\in N_D^+(w)\cup N_D^{++}(w),
\]
or
\[
 \forall w\notin\{a,b\}:\quad
 w\to b\ \Longrightarrow\ a\in N_D^+(w)\cup N_D^{++}(w).
\]
The respective convenient orientations are $a\to b$ and $b\to a$.
The following published result has no generalized-star assumption
on the missing graph.

\begin{theorem}[Ghazal {\cite[Theorem~2]{Ghazal2012}}]
\label{thm:all-missing-good}
If all missing pairs of a nonempty oriented graph are good, that graph
has a Seymour vertex.
\end{theorem}
\begin{proof}[Unweighted completion proof]
Conveniently orient every missing pair to obtain a tournament $T$.
Choose an order $L$ maximizing the number of forward arcs and let
$f$ be its last vertex.
Reverse the $r$ added arcs directed out of $f$, obtaining $T'$.
The order $L$ gains $r$ forward arcs; any other order gains at most
$r$.  Hence $L$ remains a median order.
The published feed-vertex theorem
\cite{HavetThomasse2000} gives the SNP for $f$ in $T'$.

All added arcs incident to $f$ now point inward, so
$N_{T'}^+(f)=N_D^+(f)$.
For a path $f\to y\to z$ in $T'$, the first arc is original.
The second arc is original or an untouched convenient orientation:
it cannot involve $f$, as $z=f$ would make a digon.
Consequently $z$ is reachable from $f$ within two steps in $D$.
If $z\in N_{T'}^{++}(f)$, it is not a direct outneighbor in $D$,
so $z\in N_D^{++}(f)$.
The tournament SNP therefore transfers to $D$.
\end{proof}

\subsection{A bad missing edge has two positive-residual witnesses}

\begin{lemma}\label{lem:bad-missing-positive-witnesses}
If a missing pair $ab$ is not good, there are distinct vertices
$u,v$ which are nonadjacent and satisfy
\[
 \eta(u)>0,\qquad \eta(v)>0.
\]
\end{lemma}
\begin{proof}
Failure of both convenient orientations supplies
$u,v\in V(D)\setminus\{a,b\}$ with
\[
 u\to a,\quad b\in U_D(u),\qquad
 v\to b,\quad a\in U_D(v).
\]
They are distinct, since $u=v$ would contradict $u\to a$ and
$a\in U_D(v)$. An arc $u\to v$ would let $u$ reach $b$ in
two steps, and an arc $v\to u$ would let $v$ reach $a$ in
two steps.  Thus $uv$ is a missing pair.

If $\eta(u)=0$, Lemma~\ref{lem:zero-residual-source-protection}
at its far target $b$ gives
$N_D^-(b)\subseteq N_D^-(u)\setminus\{b\}$.
But $v\to b$ would imply $v\to u$, a contradiction.
Symmetrically, $\eta(v)=0$ would imply $u\to v$.
Nonnegativity and integrality of the residuals prove the claim.
\end{proof}

\begin{theorem}[Strict missing-pair margin at all orders]
\label{thm:strict-missing-margin}
The positive-residual set $Z=\{x:\eta(x)>0\}$ induces a missing pair.
In particular,
\[
 {\sum_x\eta(x)\ge2,\qquad
 M\ge\left\lceil\frac n2\right\rceil+1.}
\]
\end{theorem}
\begin{proof}
Since $D$ is all-deficient, Theorem~\ref{thm:all-missing-good}
implies that some missing pair is not good.
Lemma~\ref{lem:bad-missing-positive-witnesses} places both endpoints
of a missing witness pair inside $Z$.
Each contributes at least one, so $2M-n=\sum_x\eta(x)\ge2$.
Integer rounding gives $M\ge\lceil(n+2)/2\rceil
=\lceil n/2\rceil+1$.
\end{proof}

\begin{corollary}\label{cor:strict-first-dense-layer}
At order $n=2\delta+3$,
\[
 {M\ge\delta+3}
\]
for every minimum outdegree $\delta$.
The one-unit residual layer $2M-n=1$ is impossible at every odd order,
and the zero-residual layer is impossible at every even order.
\end{corollary}


The strict margin above is an intermediate step; the parity-sensitive final bound is Theorem~\ref{thm:supp-residual-seven}.
